\documentclass[aspectratio=169, 11pt]{beamer}
\usepackage{beamer_preamble}

\begin{document}

% ==== Title ===============================================================
\titleslide{Lesson 4-1 \textperiodcentered\ Period 3}{Reciprocal Function + Translations --- Assessment-Critical}

% ==== Do Now ==============================================================
\begin{frame}{Do Now --- Bridge from P2}{Algebra 2 \textperiodcentered\ Lesson 4-1 \textperiodcentered\ Period 3}
\phasebadges{1}{5}{bridge from P2}
\vspace{0.2cm}
\begin{projcard}{Practice \#13 \textperiodcentered\ Generalize \(k\)}{slidegold}
For an inverse variation, write an equation that gives the value of \(k\) in terms of \(x\) and \(y\).

\vspace{0.25cm}
\textbf{Then predict:} Sketch what you think \(y = 1/x\) looks like.\\
\small (No wrong answers --- just predict. We check in Launch.)
\end{projcard}
\end{frame}

% ==== Launch ==============================================================
\begin{frame}{Launch --- Ex 4 + Ex 5 (teacher models both)}{Algebra 2 \textperiodcentered\ Lesson 4-1 \textperiodcentered\ Period 3}
\phasebadges{2}{12}{teacher models}
\vspace{0.1cm}
\begin{projcard}{Parent: \(f(x) = 1/x\)}{slidegreen}
\begin{itemize}
  \setlength{\itemsep}{0pt}
  \item V.A.: \(x = 0\) \quad H.A.: \(y = 0\)
  \item Domain: \(x \neq 0\) \quad Range: \(y \neq 0\)
  \item Hyperbola, Quadrants I and III. Never touches either axis.
\end{itemize}
\end{projcard}
\vspace{-0.2cm}
\begin{projcard}{Translated: \(g(x) = 1/(x-h) + k\)}{slidegreen}
\begin{itemize}
  \setlength{\itemsep}{0pt}
  \item V.A. \(\rightarrow x = h\); H.A. \(\rightarrow y = k\)
  \item Domain: \(x \neq h\) \quad Range: \(y \neq k\)
  \item \((h, k)\) = \textbf{new center} (where asymptotes cross)
  \item \textcolor{slideaccent}{Ex: \(g(x) = 1/(x-3)+2\) \(\rightarrow\) asymptotes \(x=3\), \(y=2\)}
\end{itemize}
\end{projcard}
\end{frame}

% ==== Explore =============================================================
\begin{frame}{Explore --- TPS: Parent Forms}{Algebra 2 \textperiodcentered\ Lesson 4-1 \textperiodcentered\ Period 3}
\phasebadges{2}{35}{student work}
\small\textcolor{slidegray}{6 items total. Parent forms first (laps 1--3), then translations (laps 4--6).}

\vspace{0.1cm}
\begin{projcard}{Try It 4 \textperiodcentered\ Scaled parent}{slideblue}
Graph \(g(x) = 10/x\). State domain, range, and asymptotes.
\end{projcard}
\vspace{-0.2cm}
\begin{projcard}{Practice \#11 \textperiodcentered\ Error Analysis: \(y = 5/x\)}{slideblue}
Describe and correct the student's graphing error.
\end{projcard}
\vspace{-0.2cm}
\begin{projcard}{Practice \#18 \textperiodcentered\ Reflected parent: \(y = -2/x\)}{slideblue}
Graph \(y = -2/x\). State domain, range, and asymptotes.
\end{projcard}
\end{frame}

% ==== Explore (translations) ===============================================
\begin{frame}{Explore --- TPS: Translations}{Algebra 2 \textperiodcentered\ Lesson 4-1 \textperiodcentered\ Period 3}
\phasebadges{2}{35}{student work (cont.)}
\small\textcolor{slidegray}{Before each translation: identify \((h,k)\), state asymptotes, then graph.}

\vspace{0.1cm}
\begin{projcard}{Try It 5 \textperiodcentered\ Translated: \(g(x) = 1/(x+2) - 4\)}{slideblue}
Graph \(g(x) = 1/(x+2) - 4\). Equations of asymptotes? Domain and range?
\end{projcard}
\vspace{-0.2cm}
\begin{projcard}{Practice \#19 \textperiodcentered\ \(\bigstar\) ASSESSMENT CRITICAL}{slideaccent}
\(g(x) = 1/(x-2) + 6\). Asymptotes, domain, range, \textbf{and intercepts}. (\(\sim\)8 min)
\end{projcard}
\vspace{-0.2cm}
\begin{projcard}{Practice \#21 \textperiodcentered\ Video game storage (if time)}{slideblue}
Number of games \(n\) varies inversely with average game size \(s\) (GB).
\end{projcard}
\end{frame}

% ==== Share / Summary =====================================================
\begin{frame}{Share / Summary --- Concept Wrap + Assessment Preview}{Algebra 2 \textperiodcentered\ Lesson 4-1 \textperiodcentered\ Period 3}
\phasebadges{2}{5}{DOK-2 reveal + LEHS preview}
\vspace{0.1cm}
\begin{projcard}{Sentence frame \textperiodcentered\ Practice \#19 pair share}{slideblue}
\small
``This graph is the parent \(f(x) = 1/x\) translated \underline{\hspace{0.8cm}} units \underline{\hspace{0.8cm}} and \underline{\hspace{0.8cm}} units \underline{\hspace{0.8cm}}.
The asymptotes are \(x =\) \underline{\hspace{0.5cm}} and \(y =\) \underline{\hspace{0.5cm}}.
The intercepts are \underline{\hspace{1.8cm}}.''
\end{projcard}

\vspace{-0.1cm}
\begin{projcard}{\(\bigstar\) Assessment Signal --- Practice \#19 is the KEYSTONE REHEARSAL}{slideaccent}
\small
\textbf{Practice \#19} is the keystone rehearsal item for Topic 4 LEHS:
\begin{itemize}
  \setlength{\itemsep}{0pt}
  \item \textbf{Q\#3} (H/V asymptotes) --- \#19 asked for exactly this.
  \item \textbf{Q\#5} (translation description) --- the sentence frame above IS the answer shape.
  \item \#19 hit BOTH skills in one item. If you got \#19, you are ready.
\end{itemize}
\end{projcard}
\end{frame}

% ==== Exit Ticket =========================================================
\begin{frame}{Exit Ticket --- Pre-Assessment Confidence Check}{Algebra 2 \textperiodcentered\ Lesson 4-1 \textperiodcentered\ Period 3}
\phasebadges{1-2}{3}{not graded}
\vspace{0.2cm}
\begin{projcard}{Complete on your exit ticket}{slidegold}
\begin{enumerate}
  \setlength{\itemsep}{0.2cm}
  \item I can identify asymptotes of \(y = 1/(x-h) + k\). \hfill Confidence: 1 / 2 / 3 / 4 / 5
  \item I can describe a translation from \(f(x) = 1/x\) to \(g(x)\). \hfill Confidence: 1 / 2 / 3 / 4 / 5
  \item One thing I want to review before Topic 4 assessment: \underline{\hspace{3.5cm}}.
\end{enumerate}
\end{projcard}
\end{frame}

\end{document}
